\section{Probability Calculi from Certain States}
\label{sec:recipe}

Corollary~\ref{cor:tensor-dichotomy} shows that linear logic admits no
grading by Definition~\ref{def:lvaluation}'s axioms that distinguishes
$a$ from $a \otimes a$. The obstruction is in the modularity axiom, which
uses the lattice operations of the event algebra. This section uses
instead a construction that needs no lattice operations: the calculus of
a logic is the set of mixtures of its certain states, its
$\{0,1\}$-valued assignments. For classical logic this is de Finetti's
characterization of coherent probability
assignments~\cite{definetti1974}, and Paris~\cite{paris2001} and
Williams~\cite{williams2012} use it for non-classical logics, where it is
known as generalized probabilism. On finite frames, and in particular on
finite Boolean algebras, it returns exactly the valuations of
Definition~\ref{def:valuation}, an instance of Paris's theorem for logics
whose certain states respect $\wedge$ and $\vee$~\cite{paris2001}.
Linear logic is not such a logic. For every formula of
multiplicative--additive linear logic over finite data types
(Definition~\ref{def:mall}), its output is contained in the
probabilistic coherence space (Definition~\ref{def:pcs}) of the formula,
and the certain elements of that space are cliques (sets of pairwise
coherent points) of Girard's coherence space
(Section~\ref{sec:recipe-linear}). At implications between data types it
returns the probabilistic coherence space. At the exponential
$!\mathrm{Bool}$ (the type of Booleans that may be used any number of
times) it returns a strict subset.

\subsection{The Construction}

\begin{definition}[presentation by certain states, mixtures\leanref{mix}]
\label{def:mix}
A \emph{presentation} is a set $E$ of events together with a set $S$ of
$\{0,1\}$-valued functions $E \to \ENNReal$, the \emph{certain states}.
Its \emph{calculus} $\mathrm{Mix}(S)$ is the set of finite convex
combinations $\sum_{i} w_i\, s_i$ with $s_i \in S$, $w_i \in \ENNReal$,
and $\sum_i w_i = 1$.
\end{definition}

\noindent
For classical logic the certain states are the two-valued models, for
intuitionistic logic on a frame they are the sharp valuations
(Theorem~\ref{thm:sharp-iff-point}), and for a type of linear logic they
are the deterministic programs of that type
(Section~\ref{sec:recipe-linear}). Two properties hold for every
presentation.

\begin{theorem}[certain elements and mixing\leanref{sharp-mem-mix-iff}\leanref{mix-mix}]
\label{thm:mix-basic}
For every presentation $(E, S)$:
(i) a $\{0,1\}$-valued $x$ lies in $\mathrm{Mix}(S)$ iff $x \in S$, and
(ii) $\mathrm{Mix}(\mathrm{Mix}(S)) = \mathrm{Mix}(S)$.
\end{theorem}

\noindent
Part (i) is the first condition of Section~\ref{sec:intro}, that the
$\{0,1\}$-valued elements of the calculus reconstruct the logic. In a
presentation the logic is given by $S$, so the condition says that mixing
creates no new $\{0,1\}$-valued elements. The proof: if
$x = \sum_i w_i s_i$ is $\{0,1\}$-valued and $w_j \ne 0$, then at each
$e$ with $x_e = 0$ every $s_i$ with $w_i \ne 0$ has $s_i(e) = 0$, and at
each $e$ with $x_e = 1$ every such $s_i$ has $s_i(e) = 1$, since
$\sum_i w_i = 1$ and $s_i(e) \le 1$. So $s_j = x$, and $x \in S$.
Part (ii) states that a convex combination of elements of
$\mathrm{Mix}(S)$ is again in $\mathrm{Mix}(S)$. For valuations on a
frame this is the closure of valuations under mixing of
Section~\ref{sec:representation}\leanref{valuation-mix}.
The second condition of Section~\ref{sec:intro}, that the classical
case returns the classical theory, is Theorem~\ref{thm:frame-recipe}
below.

The three axioms of Definition~\ref{def:valuation} are instances of one
form:

\begin{definition}[affine law\leanref{affinelaw}]
\label{def:affine-law}
For finite $E$, an \emph{affine law} is given by $a, b : E \to \ENNReal$
and $k, k' \in \ENNReal$. A function $x : E \to \ENNReal$ \emph{satisfies}
it if $\sum_e a_e x_e + k \le \sum_e b_e x_e + k'$.
\end{definition}

\noindent
Monotonicity $x_a \le x_b$, normalization ($x_\top \le 1$ and
$1 \le x_\top$, $x_\bot \le 0$), and modularity (two opposite
inequalities) are affine laws.

\begin{theorem}[completeness\leanref{mix-eq-laws}]
\label{thm:mix-eq-laws}
For finite $E$ and every set $S$ of $\{0,1\}$-valued functions on $E$,
$\mathrm{Mix}(S)$ is the set of functions satisfying every affine law
that every $s \in S$ satisfies.
\end{theorem}

\noindent
The inclusion from left to right holds because affine laws are preserved
by convex combinations\leanref{affinelaw-holds-of-mix}. The converse
separates a point outside the convex hull of $S$ by a hyperplane, using
\textsf{mathlib}'s geometric Hahn--Banach theorem, and reads the
hyperplane off as an affine law. For rational points and integer certain
states, Section~\ref{sec:hinge} gives a proof without choice by
Fourier--Motzkin elimination (elimination of variables from a system
of linear inequalities)\leanref{constr-farkas}.
Theorem~\ref{thm:mix-eq-laws} is the generalization of the three axioms:
the axioms of the calculus of a presentation are the affine laws valid
on its certain states, and nothing else. On frames these are the three
axioms of Definition~\ref{def:valuation}:

\begin{theorem}[frames\leanref{frame-recipe}]
\label{thm:frame-recipe}
For a finite poset $P$, the valuations on the frame
$\mathrm{Low}(P)$ of lower sets of $P$ are exactly $\mathrm{Mix}$ of its
sharp valuations.
\end{theorem}

\noindent
The proof combines Theorem~\ref{thm:mixture} with the fact that a sharp
valuation on $\mathrm{Low}(P)$ is a point valuation
$\delta_p$\leanref{eq-deltapoint-of-issharp}. Every finite frame is of
the form $\mathrm{Low}(P)$ (Section~\ref{sec:representation}), and finite
Boolean algebras are the case of a discrete $P$ (no two distinct
elements comparable), where the valuations are the probability mass functions
on $P$. So the construction reduces to the valuation calculus on finite
intuitionistic event algebras and to Kolmogorov's calculus on finite
classical ones.

For infinite $E$ the laws are taken with finite support. A
\emph{window} is a finite set $F \subseteq E$ of events, and a
\emph{finite law} is an affine law on some window, applied to the
restrictions of functions to $F$.

\begin{theorem}[barycenters\leanref{laws-iff-local-mix}\leanref{laws-iff-barycenter}]
\label{thm:barycenter}
Let $S$ be a set of $\{0,1\}$-valued functions on an arbitrary set $E$.
For $x : E \to \ENNReal$ the following are equivalent:
(i) $x$ satisfies every finite law valid on $S$.
(ii) For every window $F$, the restriction of $x$ to $F$ lies
in $\mathrm{Mix}$ of the restrictions of $S$ to $F$.
(iii) There is a Borel probability measure $\mu$ on $\{0,1\}^E$, with
the product topology (the topology of convergence at each event), such
that
$x_e = \mu\{\sigma \mid \sigma_e = 1\}$ for every $e \in E$, and
$\mu(\mathrm{bad}_F) = 0$ for every window $F$, where $\mathrm{bad}_F$ is
the set of $\sigma$ whose restriction to $F$ is the restriction of no
$s \in S$.
\end{theorem}

\noindent
The equivalence of (i) and (ii), which Paris states for infinite
languages~\cite{paris2001}, applies Theorem~\ref{thm:mix-eq-laws} to
each window. For (ii)~$\Rightarrow$~(iii), the mixture on a window $F$
gives a probability measure $\mu_F$ on $\{0,1\}^E$, a finite convex
combination of point masses at certain states. The windows are directed
by inclusion. The probability measures on the compact space
$\{0,1\}^E$ form a compact space in the weak topology (the
coarsest topology in which $\nu \mapsto \int f \, d\nu$ is continuous
for every bounded continuous $f$)\footnote{In \textsf{mathlib} this is an instance proved from
Prokhorov's theorem.}, so the family $(\mu_F)_F$ has a cluster point
$\mu$. A \emph{cylinder set} is a set of $\sigma$ determined by the
values of $\sigma$ on a window. Cylinder sets are clopen (closed and
open), so
$\nu \mapsto \nu(C)$ is continuous in the weak topology for every
cylinder set $C$. The sets $\{\sigma \mid \sigma_e = 1\}$ and
$\mathrm{bad}_F$ are cylinder sets, and $\mu_F$ has the required values
on them once $e \in F$, so $\mu$ has them. For (iii)~$\Rightarrow$~(i),
a finite law valid on $S$ holds at $\mu$-almost every $\sigma$, and
integrating it against $\mu$ gives the law at $x$.

\paragraph{Relation to generalized probabilism.}
Theorem~\ref{thm:mix-eq-laws} is de Finetti's theorem in the form used
for non-classical logics~\cite{paris2001,williams2012}. An affine law
valid on every certain state and violated at $x$ is a finite system of
bets priced by $x$ with a net loss in every certain state, a Dutch book,
so the theorem says that $x$ admits no Dutch book iff it is a mixture of
certain states. Paris proves it with a separating hyperplane, as the Lean
proof does, and Williams derives from the same fact that these are the
assignments not dominated in accuracy~\cite{williams2012}. Paris also
shows that when every certain state respects $\wedge$ and $\vee$ the
affine laws reduce to normalization, monotonicity, and
modularity~\cite{paris2001}, which gives Theorem~\ref{thm:frame-recipe}.
Theorem~\ref{thm:barycenter}(iii) is the counterpart, for
$\{0,1\}$-valued certain states, of the representation of the states of
an MV-algebra (an algebra of \L{}ukasiewicz many-valued logic) as
integrals over its
valuations~\cite{kroupa2006,panti2008,kuhr-mundici2007}. Gil Sanchez, Gyenis, and Wro\'nski axiomatize the convex hull for a class
of non-classical evaluations and question the use of Dutch books outside
classical logic~\cite{gilsanchez2022}, Bradley studies these convex hulls
for non-classical probability~\cite{bradley2016}, and Gyenis proves that for
every logic given by a finite algebra of truth values, each with a
numerical value, the convex hull of the evaluations has an effective
axiomatization in an equational calculus, finite for finitely many
variables and not finite in general~\cite{gyenis2025}.
Theorem~\ref{thm:mix-eq-laws} is the elementary counterpart without a
formula language: it applies to an arbitrary set of $\{0,1\}$-valued
certain states on a finite set of events, and its laws are individual
affine inequalities, not schemata uniform in the variables. This is what
lets it apply to the web of a linear formula, whose certain states are
not evaluations into a fixed algebra of truth values. This section
adds the mechanization of these results, a proof of Theorem~\ref{thm:mix-eq-laws}
without choice for rational data, and their use for logics that
Paris's axiomatization does not cover: for linear logic, where modularity is not available
(Corollary~\ref{cor:tensor-dichotomy}), in Sections~\ref{sec:recipe-linear}
and~\ref{sec:second}, and for quantum event logic in
Section~\ref{sec:quantum}.

\subsection{Uniqueness}
\label{sec:recipe-unique}

$\mathrm{Mix}(S)$ is the least convex set containing $S$. Criteria
(1) and (2) of Section~\ref{sec:intro} do not single it out: the
probabilistic coherence space $!\mathrm{Bool}$ satisfies (1) and strictly
contains the calculus (Theorem~\ref{thm:bang}). Criterion (3) does:
every affine law that holds on all certain states holds on the whole
calculus. These laws are, relative to the logic, Boole's conditions of
possible experience~\cite{pitowsky1994}.

\begin{theorem}[uniqueness\leanref{calculus-unique}]
\label{thm:unique}
For finite $E$ and a set $S$ of $\{0,1\}$-valued functions on $E$, let
$C$ contain $S$, be closed under finite mixtures, and satisfy every
affine law that every $s \in S$ satisfies. Then $C = \mathrm{Mix}(S)$.
\end{theorem}

\noindent
Closure under mixtures gives $\mathrm{Mix}(S) \subseteq C$, and
Theorem~\ref{thm:mix-eq-laws} gives $C \subseteq \mathrm{Mix}(S)$. The
sets that contain $S$ and are closed under mixtures have a least element,
$\mathrm{Mix}(S)$. The sets that satisfy the laws of $S$ have a greatest
element, the set of solutions of those laws. The theorem holds because
the two coincide, which is Theorem~\ref{thm:mix-eq-laws}. The mathematics
is de Finetti's and Paris's~\cite{definetti1974,paris2001}. What the
theorem adds is the use of the three conditions as criteria for a
calculus.

For infinitely many events one more condition is needed, since
$\mathrm{Mix}(S)$ is in general not closed and so is not the solution set
of any family of laws. A state is checked against finitely many events
at a time. We therefore require the calculus to be closed in the product
topology on $\ENNReal^{E}$, the topology of convergence at each event: a
function that is approximated by states of the calculus on every finite
window belongs to the calculus.

\begin{theorem}[uniqueness for arbitrary events\leanref{calculus-unique-closed}]
\label{thm:unique-closed}
For an arbitrary set $E$ and a set $S$ of $\{0,1\}$-valued functions on
$E$, let $C$ contain $S$, be closed under finite mixtures and in the
product topology, and satisfy every finite law that every $s \in S$
satisfies. Then $C$ is the closure of $\mathrm{Mix}(S)$, and it is the
set of functions satisfying every finite law valid on $S$.
\end{theorem}

\noindent
A function satisfying every finite law valid on $S$ lies, on each window
$F$, in $\mathrm{Mix}$ of the restrictions of $S$ to $F$
(Theorem~\ref{thm:barycenter}). So every basic neighbourhood of it,
which constrains finitely many events, meets
$\mathrm{Mix}(S)$\leanref{localmix-subset-closure}. The rest of the proof
is that of Theorem~\ref{thm:unique}.

\begin{corollary}[valuations on frames]
\label{cor:frame-unique}
In a classical metatheory, for every frame $\Omega$ the valuations on
$\Omega$ are the closure of $\mathrm{Mix}$ of its sharp valuations. So
they are the only set of functions on $\Omega$ that contains the sharp
valuations, is closed under mixtures and in the product topology, and
satisfies every finite law valid on the sharp valuations.
\end{corollary}

\noindent
A finite set of events generates a finite distributive sublattice $L$
of $\Omega$. The restriction of a valuation to $L$ is a valuation on
$L$, hence a mixture of sharp valuations on $L$
(Theorem~\ref{thm:frame-recipe}, with $L \cong \mathrm{Low}(P)$ by
Birkhoff's theorem, which presents a finite distributive lattice as
the lower sets of its join-irreducible elements). Each sharp valuation on $L$ is the restriction of
a sharp valuation on $\Omega$, by the prime ideal theorem. So every
valuation satisfies the condition of Theorem~\ref{thm:barycenter}(ii),
and the valuations satisfy every finite law valid on the sharp
valuations. Conversely, normalization, monotonicity, and modularity are
such laws. Theorem~\ref{thm:unique-closed} then applies. This
corollary is not formalized.

\paragraph{Why the third criterion.}
The third criterion has four readings, and on finite $E$ they agree.
\begin{itemize}
  \item \emph{Coherence.} A law valid on $S$ and violated at $x$ is a
  Dutch book against $x$, so the criterion excludes every state against
  which the logic guarantees a loss~\cite{paris2001}.
  \item \emph{Accuracy.} Williams shows that, for the Brier score (the sum over events of the
  squared distance to the truth value), the states
  outside $\mathrm{Mix}(S)$ are exactly those dominated in accuracy by a
  state inside it~\cite{williams2012}. The criterion excludes dominated states.
  \item \emph{Ignorance.} For finite $E$, a compact convex set of real-valued
  functions on $E$ is the convex hull of its extreme points (the points
  that are not mixtures of other points of the set), by
  Minkowski's theorem. So a compact convex
  set containing $S$ satisfies the criterion iff its extreme points are
  certain states. Uncertainty is then uncertainty about which certain
  state obtains. This equivalence is not formalized.
  \item \emph{No new behaviours.} Every known failure of the criterion is
  a state that no mixture of certain states produces: the point of
  $!\mathrm{Bool}$ that violates the law $x_{(1,1)} \le 0$, valid on its
  certain states, the points of the probabilistic coherence space at
  second order outside the calculus (Section~\ref{sec:second}), and the
  quantum states on Cabello's set, which violate the bound of
  Theorem~\ref{thm:quantum-partial}.
\end{itemize}
The criterion is therefore a choice and not a consequence of (1) and
(2). A calculus that drops it is either a semantics of programs with more
behaviours than its certain states, as for probabilistic coherence
spaces, or a contextual theory (one whose states are certain only
within each context, Section~\ref{sec:quantum-local}), as for quantum
states. How the criteria
should treat logics without enough certain states is the subject of
Section~\ref{sec:quantum-local}.

Each condition of Theorem~\ref{thm:unique} is needed. Without (3),
$!\mathrm{Bool}$ is a counterexample. Without closure under mixtures,
$S$ itself is one whenever $S$ has two elements. Without $S \subseteq C$,
every proper convex subset of $\mathrm{Mix}(S)$ is one. Criterion (2) is a
consequence of the others: at a finite Boolean algebra the certain states
are the point valuations of its atoms, and their mixtures are the
probability mass functions on the atoms (Theorem~\ref{thm:frame-recipe}).
On infinite Boolean algebras they give the finitely additive
probabilities (Corollary~\ref{cor:frame-unique}), so countable additivity
would be an additional requirement. Under (3) the criterion is redundant,
but it still constrains calculi that drop (3). The local calculus of
Section~\ref{sec:quantum-local} and the probabilistic coherence spaces
violate (3) and satisfy (2): with a single context the local calculus is
classical probability on the outcomes, and at first order the
probabilistic coherence space is the calculus (Theorem~\ref{thm:pcs-recipe}).

\paragraph{Not conflating events.}
Criteria (1) and (3) are the two halves of one condition. A set that
contains every certain state satisfies no affine law that some certain
state violates. So under (1) the calculus imposes no law that the logic
does not have. Criterion (3) is the converse. Together they say that the
calculus satisfies exactly the affine laws of the certain
states\leanref{calculus-laws-iff}. On a frame this excludes two kinds of
conflation, in a classical metatheory.
\begin{itemize}
  \item \emph{Of events.} If $a \ne b$, a sharp valuation separates them
  (Theorem~\ref{thm:sharp-separating}), so some state of the calculus
  gives $a$ and $b$ different values.
  \item \emph{Of negation with complement.} If every sharp valuation
  satisfies $v(a) + v(\neg a) = 1$, then $a \sqcup \neg a =
  \top$\leanref{em-of-sharp-compl}. So a calculus that contains the
  certain states obeys the complement rule at $a$ only where excluded
  middle holds at $a$.
\end{itemize}
The classical alternative has the same numbers and different events. A
valuation on a frame extends to a finitely additive probability on the
Boolean algebra the frame generates~\cite{pettis1951}, and reading $\neg
a$ in that algebra as the Boolean complement of $a$ assigns it $1 - v(a)$.
This grading satisfies the complement rule and gives the event ``$a$ is
not verified'' the credence of the event ``$a$ is refuted''. The second
item says that no calculus satisfying (1) on the frame does this. The
difference is the slack, which is the boundary mass in the measure model
(Theorem~\ref{thm:slack-frontier}). For linear logic the analogous
conflation is of $a$ with $a \otimes a$, which every grading of the event
lattice by Definition~\ref{def:lvaluation} makes
(Corollary~\ref{cor:tensor-dichotomy}) and the calculus of
Section~\ref{sec:recipe-linear} does not.

\subsection{Linear Logic}
\label{sec:recipe-linear}

\begin{definition}[probabilistic coherence space~\cite{danos-ehrhard2011}\leanref{ispcs-lin}]
\label{def:pcs}
For a set $W$ and a set $Q$ of functions $W \to \ENNReal$, let the
\emph{orthogonal} of $Q$ be $Q^{\perp} = \{y : W \to \ENNReal \mid \sum_{w} x_w y_w \le 1 \text{ for all }
x \in Q\}$. A \emph{probabilistic coherence space} (PCS) with \emph{web}
$W$ is a set $P$ of functions $W \to \ENNReal$ such that
$P = P^{\perp\perp}$, every $w \in W$ has some $x \in P$ with $x_w \ne 0$
(totality), and $\sup_{x \in P} x_w < \infty$ for every $w \in W$
(boundedness).
\end{definition}

\noindent
On a finite web, $P = P^{\perp\perp}$ holds iff $P$ contains $0$ and is
closed, convex, and downward closed, by Girard's bipolar
theorem~\cite{girard2004}. The map $Q \mapsto Q^{\perp}$ is the
antiblocking duality of polyhedral combinatorics~\cite{fulkerson1971},
which Section~\ref{sec:second-perfect} uses.

\noindent
We use finite webs. For a finite set $X$, the \emph{data type} $X$ is the
PCS of $x : X \to \ENNReal$ with $\sum_a x_a \le 1$. The PCS
$X \multimap Y$ on $X \times Y$ is the set of $t$ with
$\sum_b \sum_a x_a\, t(a,b) \le 1$ for every $x$ in $X$, equivalently
$\sum_b t(a,b) \le 1$ for every $a$\leanref{mem-lin-iff}. Both satisfy
the three conditions\leanref{ispcs-lin}. A \emph{coherence
space}~\cite{girard1987} is a set with a reflexive symmetric relation,
its coherence, and a \emph{clique} is a set of pairwise coherent
elements. In the coherence space $X \multimap Y$ of two data types, two
points $(a,b)$ and $(a',b')$ cohere iff they are equal or $a \ne a'$.

Probabilistic coherence spaces interpret every formula of linear logic,
including the exponentials (the modalities $!$ and $?$, which let an
argument be used any number of times), and model probabilistic PCF (a
typed functional language with recursion and probabilistic
choice)~\cite{danos-ehrhard2011}. The construction of
Definition~\ref{def:mix} applies to every such type, with the
$\{0,1\}$-valued elements of its probabilistic coherence space as
certain states. At first order, the implications between data types,
the two agree (Theorem~\ref{thm:pcs-recipe}), and the exponential $!\mathrm{Bool}$
separates them (Theorem~\ref{thm:bang}).
Section~\ref{sec:recipe-mall} proves the facts that hold at every
multiplicative--additive formula, and Section~\ref{sec:second-other}
shows that the two differ at every formula with a negative occurrence
(an occurrence under an odd number of negations) of a second-order core
(the formula $(\mathbf{n} \multimap \mathbf{m}) \otimes (\mathbf{n'}
\multimap \mathbf{m'})$ of Section~\ref{sec:second-seq}). Whether they agree at the remaining types is open.

\begin{theorem}[first order\leanref{pcs-recipe}\leanref{sharp-lin-iff-clique}]
\label{thm:pcs-recipe}
$X \multimap Y = \mathrm{Mix}\{\det f \mid f : X \to Y_\bot\}$, where
$\det f(a,b) = 1$ if $f(a) = b$ and $0$ otherwise, and $Y_\bot$ adds a
divergent value. A $\{0,1\}$-valued $t$ lies in $X \multimap Y$ iff its
support is a clique of the coherence space $X \multimap Y$.
\end{theorem}

\noindent
At first order the certain states are the deterministic partial
programs, they are the cliques of the coherence-space model, and the
calculus is the PCS.

The PCS $!\mathrm{Bool}$ has web the finite
multisets over $\{t, f\}$, and is the biorthogonal closure
$(\cdot)^{\perp\perp}$ of the
promotions $x^{!}(i,j) = x_t^{\,i} x_f^{\,j}$ of the elements $x$ of
$\mathrm{Bool}$, where $(i,j)$ is the multiset with $i$ copies of $t$ and
$j$ copies of $f$.

\begin{theorem}[exponential\leanref{sharp-bang-mixed-zero}\leanref{prom-half-not-mix}]
\label{thm:bang}
Every $\{0,1\}$-valued element of $!\mathrm{Bool}$ vanishes at every
$(i,j)$ with $i, j \ge 1$. The promotion of $x = (\tfrac12, \tfrac12)$
takes the value $\tfrac14$ at $(1,1)$, so it lies in $!\mathrm{Bool}$ and
not in $\mathrm{Mix}$ of the $\{0,1\}$-valued elements.
\end{theorem}

\noindent
Danos and Ehrhard compute the largest value of an element of
$!\mathrm{Bool}$ at $(i,j)$, which is $i^i j^j / (i+j)^{i+j}$, below $1$
when $i, j \ge 1$, and note that these multisets are not points of the web
of $!\mathrm{Bool}$ in Girard's coherence spaces~\cite{danos-ehrhard2011}.
The first statement of Theorem~\ref{thm:bang} follows, and the second is
a direct consequence of the first. The Lean proof of the first uses the
dual element $4\,\delta_{(i,j)}$, which is admissible because
$x_t + x_f \le 1$ implies $4 x_t x_f \le 1$.

\subsection{All Multiplicative--Additive Formulas}
\label{sec:recipe-mall}

\begin{definition}[formulas and their semantics\leanref{linear-fm}]
\label{def:mall}
The formulas are generated from the data types $\mathbf{n}$, for
$n \in \mathbb{N}$, by linear negation $A^{\perp}$, $A \otimes B$, and
$A \mathbin{\&} B$. Write $A \parr B = (A^{\perp} \otimes B^{\perp})^{\perp}$,
$A \oplus B = (A^{\perp} \mathbin{\&} B^{\perp})^{\perp}$, and
$A \multimap B = (A \otimes B^{\perp})^{\perp}$. The web $|A|$ is
$\{1, \dots, n\}$ for $\mathbf{n}$, $|A|$ for $A^{\perp}$, $|A| \times |B|$
for $A \otimes B$, and $|A| \sqcup |B|$ for $A \mathbin{\&} B$. The PCS
$P(A)$ is the data type for $\mathbf{n}$, $P(A)^{\perp}$ for $A^{\perp}$,
$\{x \otimes y \mid x \in P(A),\ y \in P(B)\}^{\perp\perp}$ for
$A \otimes B$, where $(x \otimes y)(a,b) = x_a y_b$, and the set of $z$
whose restrictions to $|A|$ and $|B|$ lie in $P(A)$ and $P(B)$ for
$A \mathbin{\&} B$~\cite{danos-ehrhard2011}. The coherence $\coh_A$ is
equality on $\mathbf{n}$, and $s \coh_{A^{\perp}} t$ iff $s = t$ or not
$s \coh_A t$. On $A \otimes B$ it is componentwise. On $A \mathbin{\&} B$
it is $\coh_A$ and $\coh_B$ on the two components, and any two points from
different components cohere~\cite{girard1987}. The \emph{coherence
graph} $G_A$ of $A$ has the web $|A|$ as vertices and an edge between
any two distinct coherent points.
\end{definition}

\begin{theorem}[every formula\leanref{ispcs-p}\leanref{linear-invariants}\leanref{sharp-isclique}\leanref{recipe-subset-p}]
\label{thm:mall}
For every formula $A$:
(i) $P(A)$ is a PCS.
(ii) The indicator of every point and of every pair of distinct coherent
points lies in $P(A)$, and $x_s + x_t \le 1$ for every $x \in P(A)$ and
every pair of distinct incoherent points $s, t$.
(iii) Every $\{0,1\}$-valued element of $P(A)$ is the indicator of a
clique of the coherence space $A$.
(iv) $\mathrm{Mix}$ of the $\{0,1\}$-valued elements of $P(A)$ is
contained in $P(A)$, and its $\{0,1\}$-valued elements are exactly those
of $P(A)$.
\end{theorem}

\noindent
Part (ii) is proved by simultaneous induction on $A$. Negation exchanges
the statement about coherent pairs with the statement about incoherent
pairs. For $A \otimes B$, the indicator of a coherent pair is bounded by
the tensor of indicators of its projections, and a pair incoherent in
$A \otimes B$ is incoherent in one component, where the bound of $A$ or
$B$ applies. Part (iii) follows from (ii): two incoherent points in the
support of a $\{0,1\}$-valued element would give it total $2$ on a pair
bounded by $1$. Part (iv) is Theorem~\ref{thm:mix-basic} together with
the closure of an orthogonal $Q^{\perp}$ under mixing\leanref{mix-subset-orth}.
Girard introduces probabilistic coherence spaces with the intuition that
two incoherent points of the web have weights summing to at most $1$, and then sets
coherence aside~\cite{girard2004}. Part (ii) shows that the intuition
holds at every multiplicative--additive formula over data types. The
converse of (iii) fails (Section~\ref{sec:second-shannon}).

The semantics of Definition~\ref{def:mall} computes the spaces of
Section~\ref{sec:recipe-linear}: $P(\mathbf{n} \multimap \mathbf{m})$ is
the PCS $X \multimap Y$ of that section for $X = \{1,\dots,n\}$ and
$Y = \{1,\dots,m\}$\leanref{p-lolli-base}. So
Theorem~\ref{thm:pcs-recipe} is a statement about these
formulas\leanref{recipe-formula}.

The exponentials, whose
webs are infinite, are not covered by Definition~\ref{def:mall}.
Theorem~\ref{thm:bang} treats $!\mathrm{Bool}$ separately.
